% -----------------------------------------------------------------------
% TWIN-LOCATION NOTICE. This file is one of the shared homes for the same
% results; edit every listed home or the drift checker fails.
%   Standalone paper : docs/harbor-research/tex/paper3.tex (S4-S7)
%   Chapters         : .../spawn-to-person.tex (Ch.4, S(oracle));
%                       .../agent-transactions-whitepaper.tex
%                       (Ch.6, S(claim-signaling-ic))
%   Index ids        : R7, prop:claim-signaling-ic
%   Check            : python3 scripts/harbor-research/check_library_index.py
%                      python3 scripts/harbor-research/check_propagated_corrections.py
%                      python3 scripts/harbor-research/check_citations.py
%   System document  : docs/harbor-research/LIBRARY-SYSTEM.md
% -----------------------------------------------------------------------
\documentclass[11pt,a4paper]{article}
\usepackage[margin=1in]{geometry}
\usepackage{amsmath,amssymb,amsthm}
\usepackage{xcolor}
\usepackage{booktabs}
\usepackage{longtable}
\usepackage{enumitem}
\IfFileExists{titlesec.sty}{%
  \usepackage{titlesec}%
  \titleformat{\section}{\large\bfseries\color{harborblue}}{\thesection}{0.7em}{}%
  \titleformat{\subsection}{\normalsize\bfseries}{\thesubsection}{0.6em}{}%
}{}
\usepackage[colorlinks=true,linkcolor=blue!45!black,urlcolor=blue!45!black]{hyperref}
\usepackage{tikz}
\IfFileExists{pgfplots.sty}{%
  \usepackage{pgfplots}%
  \pgfplotsset{compat=1.18}%
}{}
\usetikzlibrary{arrows.meta,positioning,fit,backgrounds,calc,shapes.geometric,shapes.symbols,decorations.pathmorphing,decorations.pathreplacing}
% --- SYSTEM ARCHITECTURE COLOR PALETTE ---
\definecolor{ptRed}{HTML}{B52227}
\definecolor{ptBlue}{HTML}{274975}
\definecolor{ptDark}{HTML}{1A1A1A}
\definecolor{ptLightGray}{HTML}{F4F4F4}
\definecolor{ptGray}{HTML}{E2E2E2}
\definecolor{ptGold}{HTML}{D4A373}
\definecolor{ptGrid}{HTML}{E5E5E5}
\definecolor{harborblue}{HTML}{274975}
\definecolor{shipred}{HTML}{B52227}
\definecolor{seagreen}{RGB}{31,110,70}
\definecolor{slate}{RGB}{71,85,105}
\definecolor{slatedark}{RGB}{30,41,59}
\definecolor{amber}{HTML}{D4A373}
\definecolor{emerald}{RGB}{16,185,129}
\newtheorem{theorem}{Theorem}
\newtheorem{claim}{Claim}
\theoremstyle{definition}
\newtheorem{change}{Change}
\newtheorem{move}{Move}
\theoremstyle{remark}
\newtheorem*{verdict}{Verdict}
\newtheorem*{adj}{Adjudication}
\setlist{itemsep=1pt,topsep=2pt}

% Shared native-figure styles for the harbor-research corpus (relation-maps,
% regime diagrams, and data plots). Vector TikZ/pgfplots only.
\tikzset{
  canvasgrid/.style={step=0.25cm, draw=ptGrid, line width=0.35pt},
  swissgrid/.style={canvasgrid},
  majorgrid/.style={step=1.00cm, draw=ptGrid!75!black, line width=0.55pt},
  snode/.style={rectangle, draw=ptDark, line width=1pt, fill=white, inner sep=4pt, minimum height=0.65cm, font=\sffamily\bfseries\scriptsize, align=center},
  darknode/.style={snode, fill=ptDark, text=white, font=\sffamily\bfseries\scriptsize\color{white}},
  mask/.style={fill=white, inner sep=1.5pt},
  dArr/.style={-{Stealth[length=2.5mm, width=1.6mm]}, ptDark, line width=1.1pt, shorten >=2pt, shorten <=2pt},
  rArr/.style={-{Stealth[length=2.5mm, width=1.6mm]}, ptRed, line width=1.3pt, shorten >=2pt, shorten <=2pt},
  bArr/.style={-{Stealth[length=2.5mm, width=1.6mm]}, ptBlue, line width=1.1pt, shorten >=2pt, shorten <=2pt},
  relnode/.style={draw,thick,minimum width=2.6cm,minimum height=0.9cm,align=center,font=\small},
  relarrow/.style={-{Stealth[length=2.5mm]},thick},
  regimebox/.style={draw,thick,align=center,font=\small,inner sep=4pt}
}
\IfFileExists{pgfplots.sty}{%
\pgfplotsset{
  harbor curve/.style={
    /pgfplots/width=0.46\textwidth,/pgfplots/height=5.6cm,
    /pgfplots/axis lines=left,
    /pgfplots/label style={font=\scriptsize},
    /pgfplots/tick label style={font=\scriptsize},
    /pgfplots/title style={font=\small,align=center,yshift=2pt},
    /pgfplots/legend style={font=\scriptsize,draw=ptDark!40,line width=0.6pt,fill=white,fill opacity=0.95,text opacity=1,inner sep=3pt},
    thick,
  },
  harbor bar/.style={
    /pgfplots/width=0.46\textwidth,/pgfplots/height=5.6cm,
    /pgfplots/ybar,/pgfplots/bar width=6pt,
    /pgfplots/axis lines=left,
    /pgfplots/label style={font=\scriptsize},
    /pgfplots/tick label style={font=\tiny},
    /pgfplots/xticklabel style={rotate=40,anchor=east,font=\tiny,inner sep=1pt},
    /pgfplots/title style={font=\small,align=center,yshift=2pt},
    /pgfplots/legend style={font=\scriptsize,draw=ptDark!40,line width=0.6pt,fill=white,fill opacity=0.95,text opacity=1,inner sep=3pt},
  },
}
}{}

\usepackage{graphicx}
\graphicspath{{../figures/}}
% Hatch fill for the regime diagram's "cheating pays" half: the two regions
% must differ in texture, not only in hue, so the figure survives greyscale
% printing (fig-r7-regime.tex).
\usetikzlibrary{patterns}
\usepackage{mdframed}
\newmdenv[backgroundcolor=blue!4,linecolor=harborblue,linewidth=0.8pt,skipabove=6pt,skipbelow=6pt]{thebox}
\newmdenv[backgroundcolor=red!4,linecolor=shipred,linewidth=0.8pt,skipabove=6pt,skipbelow=6pt]{boundary}
\newcommand{\onebreath}[1]{\begin{center}\emph{#1}\end{center}}

\title{\textbf{Reputation is Amortized Verification:}\\[3pt]
\large Inspection Games for Agent Economies\\[3pt]
\normalsize Paper 3 of the Harbor research program --- B2/R7, executed}
\author{Prepared for Erich Owens}
\date{August 2026}

\begin{document}
\maketitle

\begin{abstract}
\noindent A marketplace where AI judges grade AI work can be modeled as an inspection game. We distinguish three conditional results. A bonded judge is weakly deterred iff $\rho dB\ge G$ under a keep-gain settlement rule; confiscating the gain changes the threshold to $\rho d(G+B)\ge G$. Under a sealed uniform draw from $C$ separately bribable groups, prepaid equal-price bribes give an affine payoff: bribing all groups beats bribing none iff $G>CB$. Below that threshold, and assuming a detection probability bounded below after every history of missed detections, surviving corrupt value contracts geometrically. This does not construct an honest terminal judge or prove that different model families are independent. Finally, two specified reputation-loss models produce audit costs $\Theta(\log T)$ and $O(1)$ rather than $\Theta(T)$. For the bounded-cost schedule, the exact discrete lifetime bill is $a[nG-rv n(n-1)/2]/(dB)$, where $n=\lceil G/(rv)\rceil$; $aG^2/(2dBrv)$ is its continuous approximation. These calculations identify the enforcement and measurement assumptions a deployment would have to justify.
\end{abstract}

\noindent\emph{Express lane: inspection can deter a one-shot gain; sealed sampling changes a briber's price; and accumulated forfeitable value can fund declining audit rates. Each conclusion depends on its stated payoff, information, and detection model.}

\section{Introduction}\label{sec:intro}
\textbf{The scene.} A marketplace where AI agents do work and AI judges grade it. A judge posts a bond that is slashed if it is caught cheating. A judge can take a bribe: pass shoddy work, pocket $G$. So you audit the judges --- but auditors can be bribed too, and now it looks like watchers all the way up, each level dearer than the last. Meanwhile the platform's accountant asks a different question: a judge with ten thousand verified grades is still being audited at the same rate as a newcomer --- is that spend buying anything?

\textbf{The idea in one breath.}
\onebreath{A judge has less reason to cheat when the loss on detection exceeds the gain. A sealed draw makes prepaying several independent bribe recipients expensive. A verified record can supply additional forfeitable value, allowing fewer audits without relaxing the incentive constraint.}

\textbf{The idea, before the formalism.} Start with gain $G$, audit probability $\rho$, detection probability $d$, and slashable bond $B$. Expected punishment is $\rho dB$ under the keep-gain rule. The tower calculation adds a sampling protocol: the adversary pays each chosen group before learning which group will judge. The amortization calculation adds an explicit stake $vt$ that is lost on capture. These are different extensions of the stage game, with different hypotheses. Neither a diverse list of model names nor a long history by itself supplies the needed independence, forfeiture, or detection guarantees.

\input{fig-r7-relation.tex}

\textbf{Contributions.} Section~\ref{sec:stage} separates keep-gain and confiscation settlement conventions. Section~\ref{sec:tower} derives the sealed, prepaid bribe threshold and a conditional contraction bound, then distinguishes it from an additional self-funded budget illustration. Section~\ref{sec:amort} costs two declining audit schedules, including the discrete sum rather than only its integral. Section~\ref{sec:signaling} relates this accounting to the program's separate repeated-game models. Section~\ref{sec:conjecture} states what deployment evidence is still owed. These are model calculations, not measurements of a deployed recursive judge system.

\section{Related work, and what is actually new}\label{sec:related}
The stage game is classical. Deterrence-by-expected-penalty is Becker \cite{becker1968}; the leader--follower and simultaneous inspection formulations, including the committed-inspector solution we use, are surveyed by Avenhaus, von Stengel and Zamir \cite{avz2002}, whose treatment of inspector commitment and inspectee indifference we follow. Reputation-as-equilibrium-behavior is Kreps--Wilson \cite{krepswilson1982} and Fudenberg--Levine \cite{fudenberglevine1989}: a long-lived player's continuation value disciplines short-run temptation. Batched verification of many items with few tests goes back to Dorfman's group-screening design \cite{dorfman1943}, which is the right lineage for pooled audit sampling, though we do not use pooling here. The recursive audit hierarchy of \S\ref{sec:tower} --- who audits the auditor --- is exactly the research question of Kofman and Lawarrée's collusion-in-hierarchical-agency model \cite{kofmanlawarree1993} and Tirole's three-tier principal/supervisor/agent analysis \cite{tirole1986}; the applied cryptoeconomics literature converged on the same shape from the other direction, solving the \emph{verifier's dilemma} with forced-error injection and bonded challenge-response games (TrueBit \cite{truebit2017}) and with recursive, bonded, bisecting dispute resolution (Arbitrum's interactive fraud proofs \cite{arbitrum-fraud-proofs}).

\emph{Scope of the contribution.} The stage-game solution concepts and continuation-value account are imported from this literature. The applied contribution is an explicit comparison of audit schedules and a sealed-sampling payoff calculation, with their failure conditions exposed. The equations do not establish a general solution to collusion, an implementation theorem for arbitrary LLM judges, or novelty of inspection-game deterrence itself.

\section{The vocabulary, defined}\label{sec:vocab}

\emph{Express lane: skip to \S\ref{sec:stage} --- nothing in this section is needed to follow the theorems, only to read them comfortably; the two terms that are hypotheses rather than vocabulary, \emph{sealed} and \emph{clique}, are defined where they first do work, at the head of \S\ref{sec:tower}.} This paper
is economics, and its terms are cheap to define but expensive to guess at. A software reader who has never taken
a mechanism-design course loses nothing but the dictionary; here it is, in the order the terms are needed.

An \emph{inspection game} is a two-player game between someone who may cheat and someone who may check, where checking
costs something and catching pays something. Its defining feature --- and the reason it is not simply a matter of ``audit
more'' --- is that it typically has no equilibrium in which either player behaves predictably: if the inspector always
checks, cheating never pays and the checks are wasted; if the inspector never checks, cheating always pays. The stable
outcome is that both randomise, which is why every audit rate in this paper is a probability rather than a schedule.

A \emph{mixed strategy} is exactly that randomisation --- a probability distribution over what to do rather than a choice.
Two players are in a \emph{mixed equilibrium} when each is \emph{indifferent} between its options given what the other is
doing, since a player who strictly preferred one option would simply take it. Solving for indifference is how the critical
rates below are derived, and it is the whole technique.

A player's behaviour is \emph{incentive-compatible} when honesty is at least as profitable as the alternative --- not when
it is enforced or morally preferred. Everything this paper deters, it deters by making cheating cost more than it pays.

A \emph{bond} is capital the judge posts before working and forfeits if caught cheating. It is the paper's central
instrument, and it differs from a wage in a way that matters: a wage is paid for output and withheld for non-performance
automatically, whereas a bond is forfeited only when a \emph{detection event} occurs. That gap --- between misbehaving and
being caught misbehaving --- is what the audit probability $\rho$ and detection probability $d$ measure, and it is why
$\rho d B$ rather than $B$ is the quantity a rational judge weighs. To \emph{slash} a bond is to confiscate it on
detection; the term is borrowed from proof-of-stake systems, where it means the same thing.

The two \emph{settlement conventions} of \S\ref{sec:stage} are the two natural answers to what happens to the corrupt gain
itself when a cheat is caught. Under \emph{keep-gain} the cheater keeps $G$ and loses the bond; under \emph{confiscation}
the cheater loses both. They are not a modelling detail --- they move the required audit budget by roughly 17\% at the
worked parameters, which is the paper's first practical finding.

A \emph{bribe floor} is the least a briber must offer a judge to make corruption worth its while: exactly the judge's
expected forfeiture $\beta = \rho d B$. Below that, no rational judge accepts.

Finally, \emph{amortization} is the accounting sense of the word, not the algorithmic one: spreading a fixed cost over many
uses so the per-use share falls. The paper's title claim is that a reputation \emph{is} this operation applied to
verification cost --- the reason a long-trusted judge can be audited less is that its history has already paid for the
confidence.

\section{The stage game and the convention fork}\label{sec:stage}
A judge grades one outcome. Cheating yields a one-shot gain $G$. The platform audits with probability $\rho$ at cost $a$ per audit; an audit detects a cheat with probability $d$; detection slashes the judge's posted bond $B$. Under the \emph{keep-gain} convention the caught cheat retains $G$ and forfeits $B$; under the \emph{confiscation} convention capture forfeits the gain as well.

\begin{thebox}
\begin{theorem}[Stage deterrence]\label{thm:stage}
A judge with one-shot corrupt gain $G$, facing audit probability $\rho$, detection probability $d$, and slashable bond $B$, is weakly deterred iff $\rho d B \ge G$; the critical committed audit rate is
\[ \rho^\star \;=\; \frac{G}{dB} \qquad \text{(keep-gain; Becker's condition).} \]
If capture also confiscates the gain, the cheat payoff is $G - \rho d (G{+}B)$ and the critical rate falls to
\[ \rho^\star_c \;=\; \frac{G}{d(G+B)} \;<\; \rho^\star. \]
Assume $G,B,d,a,L>0$, $d\le1$, and $G\le dB$ for a feasible committed keep-gain rate. In the simultaneous (uncommitted) game, with $a$ the audit cost and $L$ the auditor's loss from an uncaught cheat, an interior mixed Nash equilibrium exists when $G<dB$ and $a<dL$ and is $\rho^\star = G/(dB)$, $q^\star = a/(dL)$: the judge's indifference pins the audit rate, the auditor's indifference pins the cheat rate.
\end{theorem}
\end{thebox}

\emph{Why the same $\rho^\star$ twice.} The formula is unchanged across the two game forms not by coincidence but because both solve the same judge-indifference equation: commitment picks the cheapest rate that deters, and simultaneous play lands an equilibrium mix exactly there.

\emph{Proof.} Keep-gain: the cheat payoff is $G - \rho d B$, nonpositive iff $\rho \ge G/(dB)$. Confiscation: replace $B$ by $G+B$ in the forfeiture. Mixed equilibrium: the judge mixes only if indifferent, $G - \rho d B = 0$; the auditor mixes only if indifferent between auditing ($-a - qL(1-d)$) and not ($-qL$), which solves to $q^\star = a/(dL)$. Direct best-response checks confirm both indifferences numerically (\S\ref{sec:repro}). \hfill$\square$

\textbf{Numbers by hand.} At the program's running parameters $G{=}10$, $d{=}0.8$, $B{=}50$: $\rho^\star = 10/(0.8\cdot 50) = 0.25$ [verified] --- audit one grade in four, and the cheat payoff is exactly $+0.000000$ at $\rho^\star$ [internal, \texttt{b2\_\allowbreak tower.py}]. Under confiscation, $\rho^\star_c = 10/(0.8\cdot 60) = 0.2083$ [verified]. The convention moves the required audit budget from $0.25$ to $0.2083$ --- a $16.7\%$ cut for free, \emph{if and only if} your settlement rail actually claws back the gain. With $a{=}1$, $L{=}40$: $q^\star = 1/(0.8\cdot 40) = 0.03125$ [verified]. Figure~\ref{fig:regime} plots the frontier $\rho d B = G$ in the $(B,\rho)$ plane with this point marked. \emph{Now you try:} double the bond to $B{=}100$ --- what is $\rho^\star$? ($0.125$: a bigger bond lets you audit half as often; capital and vigilance are substitutes along $\rho d B = G$.)

\input{fig-r7-regime.tex}

\section{The tower: bribing sealed juries from $C$ cliques}\label{sec:tower}
A \emph{sealed} draw is hidden until the briber has committed and paid its bribes. A \emph{clique} here is a separately bribable group, not a graph-theoretic complete subgraph. Buying one group must not purchase influence over another. Distinct vendors or model families may help diversify a pool, but do not prove this economic separation or statistical independence.

\begin{thebox}
\begin{theorem}[Sealed bribe threshold and conditional contraction]\label{thm:tower}
Let an audit select one of $C$ groups uniformly, after an adversary has prepaid $\beta=\rho dB>0$ for each group it chooses to buy. Every purchased group suppresses detection; each unpurchased group audits and detects with probability $\rho d$, where $0<\rho d<1$. All bribes are paid regardless of the draw. For protected value $G_k$, the expected one-level net value from buying $c$ groups is
\[
 V(c)=G_k\left(1-\frac{C-c}{C}\rho d\right)-c\beta
     =V(0)+c\left(\frac{G_k\rho d}{C}-\beta\right).
\]
A maximum is attained at $c=0$ or $c=C$; an interior choice cannot strictly improve on both endpoints. Full bribery beats no bribery iff $G_k>CB$. Assume ties are resolved against bribery.

For a tower starting at $G_0\le CB$, additionally suppose that the probability of detection at each level, conditional on every preceding level having missed, is at least $\rho d$. Then expected surviving corrupt value obeys
\[G_k\le (1-\rho d)^kG_0.\]
For target $u>0$, no audit is needed if $G_0<u$. Otherwise
\[
 n=\left\lfloor\frac{\log(G_0/u)}{\log(1/(1-\rho d))}\right\rfloor+1
\]
levels suffice for the strict bound $G_n<u$. Bonding one selected auditor per level costs $nB$ in concurrent selected-auditor capital; it does not count collateral required from the whole candidate pool. A trustworthy terminal adjudicator remains a separate premise.
\end{theorem}
\end{thebox}

\emph{Proof.} Subtract $V(0)$ from $V(c)$. The resulting affine function has an endpoint maximum. Since $\rho d>0$, its slope is positive exactly when $G_k>CB$. With no profitable bribery and conditional detection at least $\rho d$, the probability of surviving $k$ levels is at most $(1-\rho d)^k$ by the chain rule for conditional probabilities. Multiplication by $G_0$ gives the expected-value bound. Solving $(1-\rho d)^nG_0<u$ gives the stated integer depth. A marginal detection rate at each level would not suffice: the same undetectable fault could survive every level. \hfill$\square$

\textbf{Why sealing matters.} If the adversary learns the draw before paying, it can buy only the selected group. The one-level comparison is then $G_k\rho d>\beta$, or $G_k>B$. A verifiable draw is not automatically secret before commitment: the timing and information available to the briber are part of the protocol.

\textbf{An additional budget model.} Above $CB$, the one-level theorem does not imply that corrupt value shrinks by the bribe bill. For illustration only, assume the adversary has a non-replenishing budget, this same budget determines the value it will protect next, and buying all groups reduces it by $C\beta$. This additional recursion is
\[
G_{k+1}=\begin{cases}G_k-C\beta,&G_k>CB,\\
(1-\rho d)G_k,&G_k\le CB.\end{cases}
\]
The linear prefix is a model of a self-funded budget, not a proof that paying for protection removes the underlying fraud. External financing, cross-level contracts, or correlated failures can invalidate it.

\textbf{Numbers by hand.} Let $G_0=400$, $B=50$, $\rho=0.25$, and $d=0.8$, so $\beta=10$ and $\rho d=0.2$. With $C=8$, $G_0=CB$, and ties resolved against bribery, the geometric bound applies immediately: $400,320,256,204.8,\ldots$. It takes $27$ levels to fall strictly below $u=1$, with selected-auditor bond total $1350$. Under the additional budget recursion, $C=1$ first pays $10$ per level for $35$ levels, reaches $50$, then needs $18$ geometric levels: $53$ levels and $2650$ of selected-auditor bonds. The latter comparison requires the extra budget premise. \emph{Now you try:} with $C=2$, the budget decreases by $20$ for $15$ levels to $100$; another $21$ geometric levels give $36$ in total. These are hand-recomputable model values, not measured audit performance.

\section{Reputation is amortized verification}\label{sec:amort}
The stage game prices one grade. A platform judges $T$ of them per judge lifetime, and flat auditing at $\rho^\star$ costs $a\rho^\star T = \Theta(T)$. But an old judge is not a newcomer: it has a verified history of length $t$, and with per-period reputation rent $v$ its accumulated stake $vt$ --- the future income its record commands --- is forfeit on capture. That stake is punishment capital the platform did not fund, and an incentive-compatible audit schedule may spend it.

\begin{thebox}
\begin{theorem}[Amortization]\label{thm:amort}
Assume $a,G,B,v>0$, $0<d,r\le1$, and $G\le dB$. Index grades by integers $t=0,1,\ldots,T-1$. A history of length $t$ has forfeitable monetary value $vt$; this linear value is a model assumption, not a consequence of counting grades. Let $\rho_t$ make the specified one-shot deviation weakly unprofitable.
\begin{enumerate}[leftmargin=1.6em,itemsep=1pt]
\item \textbf{Flat.} The history-blind schedule $\rho_t=G/(dB)$ spends $aGT/(dB)=\Theta(T)$.
\item \textbf{Model A: loss only if audited.} For deviation payoff $G-\rho_t d(B+vt)$, the binding schedule is $\rho_t=G/[d(B+vt)]$. Its discrete spend satisfies
\[
 \frac{aG}{dv}\log(1+vT/B)
 \le a\sum_{t=0}^{T-1}\rho_t
 \le\frac{aG}{dv}\log(1+vT/B)+\frac{aG}{dB},
\]
so it is $\Theta(\log T)$.
\item \textbf{Model B: separately specified revelation loss.} Assume deviation payoff $G-\rho_t dB-rvt$: the audit loss is the bond and the non-audit revelation loss is the reputation stake. This additive expected-loss contract must avoid counting the same forfeiture twice. Then
\[
 \rho_t=\max\{0,(G-rvt)/(dB)\}.
\]
Writing $n=\lceil G/(rv)\rceil$, the first zero-rate integer grade is $n$ and the exact lifetime spend is
\[
 S_\infty=\frac{a}{dB}\left[nG-\frac{rv\,n(n-1)}2\right]=O(1).
\]
The integral $I_\infty=aG^2/(2dBrv)$ is a continuous approximation, with $0\le S_\infty-I_\infty\le aG/(dB)$.
\end{enumerate}
All schedules lie in $[0,1]$ under the feasibility assumption and make their specified deviation payoff nonpositive. Strict deterrence requires a strict payoff inequality or a tie rule favoring honesty.
\end{theorem}
\end{thebox}

\emph{Proof.} Substitute each schedule into its deviation payoff. For Model A, compare the decreasing function $aG/[d(B+vt)]$ with its left-endpoint Riemann sum; the error lies between zero and the initial term. For Model B, only $t=0,\ldots,n-1$ have positive rate. Sum the arithmetic progression to obtain $S_\infty$. Integrating the same nonnegative decreasing ramp yields $I_\infty$ and the same initial-term error bound. The two formulae are not equal in general. \hfill$\square$

\textbf{Numbers by hand.} At $G=10,d=0.8,B=50,a=1,v=0.6,r=0.05$, the $T=200$ bills are $50$ (flat), $25.583785$ (Model A), and $35.075$ (Model B). The Model A integral is $25.495321$. Model B's real-valued zero crossing is $333\tfrac13$, so its first zero-rate integer grade is $334$. The exact lifetime bill is $41.79175$, compared with integral $41.666667$; their difference is $0.125083\ldots$, not exactly half the initial rate. At $T=1500$, flat spends $375$, Model A $61.461149$, and Model B remains at $41.79175$. These values follow directly from the displayed sums. Figure~\ref{fig:amort} plots those discrete bills.

\input{fig-paper3-amortization.tex}

\textbf{The fork is empirical.} Whether a deployment lives in Model A or Model B is a question about the world, not the model: \emph{do cheats surface without audits?} If bad grades eventually reveal themselves --- failed downstream builds, provenance tombstones, user reports --- then $r > 0$ and lifetime verification spend per honest judge is a constant. If a cheat undetected at audit time is silent forever, only the $\Theta(\log T)$ schedule is safe. The deployment telemetry that estimates $r$ (the rate at which sanctioned outcomes were first flagged by non-audit channels) is therefore part of this paper's measurement obligation, not an implementation detail.

\section{The discount-factor connection: claim signaling, mechanized}\label{sec:signaling}
The program's repeated-game track prices the same deterrence with a different currency. In the corrected claim-signaling stage game --- payoffs $(3,3)$ mutual-truth, $(0,4)/(4,0)$ sucker/defector, $(1,1)$ mutual-claim --- the one-shot deviation gain is $g = 1$ and the per-round punishment loss is $2$. By the one-shot deviation principle, a grim trigger is subgame-perfect --- cooperation survives every history, not just the equilibrium path --- iff $1 \le 2\,\delta/(1-\delta)$, i.e.\ $\delta \ge 1/3$; the gentler three-round graduated trigger is subgame-perfect iff $1 \le 2(\delta + \delta^2 + \delta^3)$, whose threshold is the unique real root in $(0,1)$ of $2\delta^3 + 2\delta^2 + 2\delta - 1 = 0$, numerically $\delta^\star = 0.3425080314$. These are not folklore numbers: they are CI-mechanized in the repository's \texttt{proofs/\allowbreak economics/\allowbreak } directory --- \texttt{delta-threshold.z3} proves existence, uniqueness in $(0,1)$, and location in $[0.34, 0.35]$ of the root by real-algebraic decision procedure; \texttt{claim\_\allowbreak signaling.tla} model-checks the discounted one-shot-deviation invariant under both TLC and Apalache; and \texttt{sweep-delta.sh} confirms the empirical crossover ($\delta = 0.35$ on the integer grid) matches the closed form. Every claim in this paragraph is [verified] by that CI suite.

The connection to Theorem~\ref{thm:amort} is the role of forfeitable continuation value. In the repeated game, a deviation is observed through that game's monitoring rule and changes future play. In Model A, an audit triggers the loss. Setting $B=0$ and $\rho=0$ in Model A removes detection and its penalty; it does not recover the repeated-game theorem. Transferring the intuition requires transferring a valid observation and punishment mechanism as well.

\section{Conjecture III.11.1, parameterized}\label{sec:conjecture}
The grading-oracle question has a conditional answer, not a general closure. A deployment must supply the following premises before the tower calculation applies.
\begin{itemize}[leftmargin=1.6em,itemsep=1pt]
\item \emph{Separately bribable groups.} Diverse models and principals are design candidates. They do not prove the absence of shared incentives, shared blind spots, or cross-group contracts.
\item \emph{A sealed uniform draw.} A verifiable random function \cite{mrv1999} can support an auditable draw. Secrecy until bribe commitment, key custody, unbiased group weighting, and resistance to recognizable honeypots need their own protocol evidence.
\item \emph{Conditional detection and a sound root.} A lower bound after histories of missed detections is stronger than an average detection score. The endpoint must be trustworthy independently of the judges it reviews; a bonded arbitration market does not become honest merely because it is named as the root.
\end{itemize}
Under these premises, $G_0\le CB$ and the tie rule yield the geometric depth bound. The $27$-level worked example prices that model. It neither supplies the premises nor establishes an implemented recursive grader.

\section{Reproducibility}\label{sec:repro}
The accompanying \texttt{b2\_\allowbreak tower.py} script lives in
\texttt{skills/\allowbreak harbor-results/\allowbreak scripts/} and uses NumPy
and Matplotlib. The program seed convention is 20260816; these calculations
are deterministic. Its output at the running parameters is reproduced below,
with long lines wrapped:
\begin{quote}\small\ttfamily
rho* (keep-gain)~~~~~= G/(dB)~~~~~~= 0.2500 \\
rho* (confiscation)~~= G/(d(G+B))~~= 0.2083 \\
q*~~~(mixed NE)~~~~~~= a/(dL)~~~~~~= 0.0312 \\
judge indifferent at rho*: cheat payoff = +0.000000 \\
C=1:~~400.0~~390.0~~380.0~~370.0~~360.0~~350.0~~340.0~~330.0~~320.0 \\
C=8:~~400.0~~320.0~~256.0~~204.8~~163.8~~131.1~~104.9~~~83.9~~~67.1 \\
levels to drive corrupt value below 1 unit (geometric regime): 27; \\
total bond capital = B per level * levels = 1350 \\
flat:~~~~~~~50.00~~~(Theta(T)) \\
model A:~~~~25.58~~~(Theta(log T);~closed form 25.50) \\
model B:~~~~35.08~~~(O(1); \\
~~closed form aG\^{}2/(2dBrv) = 41.67; hits 0 at t*=333) \\
IC check A: max deviation value = +1.78e-15 (<=0) \\
IC check B: max deviation value = +1.78e-15 (<=0)
\end{quote}
The displayed script transcript retains its own labels: its Model B ``closed form 41.67 is the continuous integral, and ``hits 0 at 333 rounds the real crossing; the integer first zero is 334 and the discrete plateau is 41.79175. Those labels must not be read as the exact discrete theorem. The script also verifies the mixed equilibrium by direct best-response check (auditor indifference $-1.25 = -1.25$) and renders the three-panel summary figure. The claim-signaling thresholds of \S\ref{sec:signaling} regenerate independently via the CI suite in \texttt{proofs/\allowbreak economics/\allowbreak } (Z3, TLC, Apalache, and the $\delta$-sweep).

\section{Honest boundaries}\label{sec:boundaries}
\begin{boundary}
\textbf{What this paper does not say.}
\begin{itemize}[leftmargin=1.6em,itemsep=2pt]
\item \emph{The headline is convention-dependent.} $\rho^\star = 0.25$ under keep-gain, $0.2083$ under confiscation --- a $\sim$17\% swing from the settlement rule alone. Quote neither without stating which convention your slashing contract implements; a platform whose escrow claws back the gain but which budgets audits at $G/(dB)$ is over-auditing, and one that budgets at $G/(d(G{+}B))$ without clawback is under-deterring.
\item \emph{Sealed sampling is a hypothesis, not a garnish.} A briber who learns the audit draw before placing bribes faces threshold $G_k > B$, not $G_k > CB$: the clique math collapses to $C{=}1$ (\S\ref{sec:tower}). The sampler --- and its VRF, and whoever operates it --- is inside the trusted computing base of Theorem~\ref{thm:tower}.
\item \emph{The tower needs an honest root.} Contraction shrinks corruption level over level; something exogenously honest must sit at the top --- an independently competent operator or adjudicator whose honesty and incentives must themselves be justified. The theorem prices the depth to that root; it does not conjure the root.
\item \emph{Detection, revelation, stake value, and cross-clique dependence are assumptions that require evidence.} Detection probability comes from audit outcomes; the revelation rate $r$ from non-audit provenance of sanctions; cross-clique collusion correlation from judge-pair agreement statistics. ``Disjoint cliques'' is falsifiable and must be monitored: two nominally rival benches fine-tuned from the same base model may fail disjointness in exactly the correlated-error cases that matter.
\item \emph{The $\Theta(\log T)$-vs-$O(1)$ fork is empirical.} Model B's constant lifetime spend holds only if cheats surface without audits ($r>0$). And at the session parameters the $O(1)$ claim is demonstrated pre-asymptotically: $n=334>T=200$, so the discrete spend $35.075$ sits below the $41.79175$ discrete lifetime bill (\S\ref{sec:amort}).
\item \emph{One-shot gains, risk-neutral bribers, stationary parameters.} $G$ is per-grade; an adversary who can correlate many small cheats below audit granularity, or who values capture asymmetrically, needs the repeated-game extension of \S\ref{sec:signaling} with the stake re-derived. The machine-epsilon IC residuals ($+1.78\times 10^{-15}$) are zeros of a float computation, stated as ``$\le 0$ up to float epsilon,'' not as exact algebra.
\end{itemize}
\end{boundary}

\begin{thebibliography}{9}
\bibitem{avz2002} R.~Avenhaus, B.~von Stengel, and S.~Zamir. Inspection games. In \emph{Handbook of Game Theory with Economic Applications}, Vol.~3, Ch.~51, pp.~1947--1987. Elsevier, 2002.
\bibitem{becker1968} G.~S. Becker. Crime and punishment: an economic approach. \emph{Journal of Political Economy} 76(2):169--217, 1968.
\bibitem{krepswilson1982} D.~M. Kreps and R.~Wilson. Reputation and imperfect information. \emph{Journal of Economic Theory} 27(2):253--279, 1982.
\bibitem{fudenberglevine1989} D.~Fudenberg and D.~K. Levine. Reputation and equilibrium selection in games with a patient player. \emph{Econometrica} 57(4):759--778, 1989.
\bibitem{dorfman1943} R.~Dorfman. The detection of defective members of large populations. \emph{Annals of Mathematical Statistics} 14(4):436--440, 1943.
\bibitem{mrv1999} S.~Micali, M.~Rabin, and S.~Vadhan. Verifiable random functions. In \emph{Proceedings of the 40th Annual Symposium on Foundations of Computer Science (FOCS)}, pp.~120--130. IEEE, 1999.
\bibitem{kofmanlawarree1993} F.~Kofman and J.~Lawarrée. Collusion in hierarchical agency. \emph{Econometrica} 61(3):629--656, 1993.
\bibitem{tirole1986} J.~Tirole. Hierarchies and bureaucracies: on the role of collusion in organizations. \emph{Journal of Law, Economics, \& Organization} 2(2):181--214, 1986.
\bibitem{truebit2017} J.~Teutsch and C.~Reitwiessner. A scalable verification solution for blockchains. arXiv:1908.04756, 2017.
\bibitem{arbitrum-fraud-proofs} Offchain Labs. Interactive fraud proofs (Arbitrum bisection dispute protocol). Technical documentation, ongoing.
\end{thebibliography}

\end{document}
